\section{What Is Portability?}
\label{sec::portability}

Two main portability types can be meant if we are talking about portability. 
On the one side, binary portability means that I compile my code once, and it runs everywhere. 
On the other side, source code portability means I write portable source code that needs to be recompiled on each new hardware platform. 
Generally, the former is much more challenging to achieve than the latter, and most of the time, some form of \gls{jit} compilation is involved. 
For the remainder of this thesis, if I talk about portability, I mean source code portability. 

The need for portable codes, in general, has been an ongoing problem for many years. 
One of the first portability definitions was already provided in \citeyear{poole_portability_1975} fifty years ago: 
\begin{quotebox}{\citet{poole_portability_1975}}
    Portability is a measure of the ease with which a program can be transferred from one environment to another: If the effort required to move the program is much less than that required to implement it initially, then we say that it is highly portable.
\end{quotebox}
% which was later picked up or slightly modified by other publications~\cite{tanenbaum_guidelines_1978, mooney_bringing_1997}. \todo{cite? more?}

\todo{find more definitions?}

Based on the slight modification of this portability definition from \citet{mooney_bringing_1997}, \citet{sedova_high-performance_2018} devised a formal definition of portability:
\begin{definition}[Portability~\cite{sedova_high-performance_2018}]
An application can be said to be portable if the cost of porting is less than the cost of re-writing:
$$DP = 1 - \left(\frac{C_P}{C_R}\right)$$
where $C_P$ is the cost to port, and $C_R$ is the cost to rewrite the program. 
\end{definition}
A program is completely portable if $DP$ is one, and a positive value indicates that porting a code is more profitable than rewriting it from scratch. 

A more recent definition of portability was created by \citet{pennycook_metric_2016, pennycook_implications_2019} which define portability as: 
\begin{quotebox}{\citet{pennycook_implications_2019}}
    The ability of an application to execute a problem correctly on a given set of platforms.
\end{quotebox}
In contrast to the previous definition, the key part here is that it is not enough for an application to run on different hardware platforms but to produce the correct results on these different platforms. 
An application that runs everywhere but produces wrong results is pointless in a real-world setting. 

\todo{insert ways to fix the portability problem}

\newpage
\subsection{OpenMP -- Open Multi-Processing}

OpenMP (Target Offloading)
\begin{itemize}
    \item first OpenMP specification version 1.0 (C/\cpp) October 1998~\cite{openmp_architecture_review_board_openmp_1998}
    \item first OpenMP specification with target offloading version 4.0 July 2013~\cite{openmp_architecture_review_board_openmp_2013}
    \item current OpenMP specification version 6.0 November 2024~\cite{openmp_architecture_review_board_openmp_2024}
    \item Implementations (Auszug)(\url{https://www.openmp.org/resources/openmp-compilers-tools/})
    \begin{itemize}
        \item AMD
        \item Intel
        \item LLVM
        \item GNU
        \item ARM
        \item NVIDIA
        \item PGI
        \item $\dots$
    \end{itemize}
\end{itemize}


\begin{itemize}
    \item \acrfull{openmp}
    \item pragma based
    \item least invasive; pragmas ignored if no OpenMP capable compiler could be found
    \item but comes also with a runtime library 
    \item originally meant for CPUs only
    \item standardized
    \item uses the fork-join model; initial thread spawns worker threads
    \item uses a custom OpemMP thread pool
    \item task- and loop-parallelism (only latter is used)
\end{itemize}


\begin{advantagebox}
    \begin{itemize}
        \item standardized
        \item supports a wide variety of hardware platforms
        \item non-intrusive and easy-to-use
    \end{itemize}
\end{advantagebox}

\begin{disadvantagebox}
    \begin{itemize}
        \item very high-level, i.e., some optimizations not possible
    \end{itemize}
\end{disadvantagebox}

\newpage
\subsection{CUDA -- Compute Unified Device Architecture}

CUDA
\begin{itemize}
    \item first release of CUDA 1.0 June, 2007~\cite{nvidia_corporation_cuda_2007}
    \item current release of CUDA 12.8 February, 2025~\cite{nvidia_corporation_cuda_2025-1}
    \item the current CUDA programming guide February, 2025~\cite{nvidia_corporation_cuda_2025}
    \item first CUDA version (3.0) deprecating device emulation mode March, 2010~\cite{nvidia_corporation_cuda_2010}
\end{itemize}

\begin{itemize}
    \item \acrfull{cuda}
    \item NVIDIA's own vendor-specific framework
    \item C-like API with custom extensions \code{<<<...>>>}
    \item only for NVIDIA \glspl{gpu}
    \item offload compute kernels from host (i.e., classically \glspl{cpu}) to one or more devices
    \item many vendor-provided, highly optimized libraries available (e.g., cuBLAS, cuSPARSE, cuFFT, $\dots$)
    \item interfaces for C, C++, FORTRAN, Python, Matlab, Java
\end{itemize}

\begin{advantagebox}
    \begin{itemize}
        \item practically, the go-to standard to target \glspl{gpu}
        \item many resources available
        \item high optimization potential
    \end{itemize}
\end{advantagebox}

\begin{disadvantagebox}
    \begin{itemize}
        \item not standardized
        \item only NVIDIA \glspl{gpu}
        \item closed-source proprietary
    \end{itemize}
\end{disadvantagebox}

\newpage
\subsection{HPX -- High Performance ParallelX}
\label{sec:frameworks:hpx}

\begin{quotebox}{\href{https://hpx-docs.stellar-group.org/latest/html/index.html\#what-is-hpx}{HPX Documentation}}
    HPX is a C++ Standard Library for Concurrency and Parallelism. It implements all of the corresponding facilities as defined by the C++ Standard. Additionally, in HPX we implement functionalities proposed as part of the ongoing C++ standardization process. We also extend the C++ Standard APIs to the distributed case. HPX is developed by the STE||AR group.
\end{quotebox}

\acrfull{hpx}~\cite{kaiser_hpx_2020} is high-performance many-task asynchronous tasking library. 
Its development started in 2007~\cite{the_stear_group_history_nodate} with the current version being v1.10.0 from March 2024~\cite{kaiser_stellar-grouphpx_2024}.
Although \gls{hpx} is not standardized where possible, its \gls{api} is closely modeled after the \cpp standard \gls{api}. 

The execution model of \gls{hpx} is based on a task graph that is implicitly built in the source code and later executed asynchronously. 
It implements custom schedulers that also incorporate work-stealing to automatically improve the load-balancing of applications using \gls{hpx}. 
This is especially useful for non-uniform problems that rely on good load balancing. 
\gls{hpx} also supports distributed memory architectures using an \gls{agas} and different accelerators using executors. 
Note that \gls{hpx} only portably incorporates kernel launches in its internal task graph; the actual compute kernels must still be written using other frameworks like \gls{cuda}, \gls{hip}, Kokkos, or SYCL.  
Additionally, \gls{hpx} supports the \cpp{17} standard library parallel algorithms (also see~\autoref{sec:frameworks:stdpar}. 
In the remainder of this thesis, only these standard algorithms are used with \gls{hpx}. 

\todo{mention that parallel algorithms were only implemented after the \cpp draft}

\begin{listing}
    \begin{minted}{cpp}
std::vector<int> indices(N);
std::iota(indices.begin(), indices.end(), 0);

hpx::for_each(hpx::execution::par_unseq, indices.begin(), indices.end(), [&](const int idx) {
    Y[idx] = alpha * X[idx] + Y[idx];
});
    \end{minted}
    \caption{%
    A simple DAXPY example using \gls{hpx} parallel loops.}
    \label{code:hpx_example}
\end{listing}

An example of these \cpp{17} standard parallel algorithms can be seen in \autoref{code:hpx_example}. 
The only difference compared to \autoref{code:stdpar_example} is the usage of the \code{hpx::} namespace qualifiers instead of the \code{std::} ones. 
In this code, the \code{for_each} loop is automatically parallelized across all available \gls{cpu} cores. 


%%% ADVANTAGES AND DISADVANTAGES
\AdvantagesAndDisadvantages{hpx}{
    \item supports algorithms close to \cpp standard
    \item task-based abstraction more suitable for some problems
}{
    \item not standardized
    \item can only launch kernels, other frameworks still \textbf{necessary} to write the actual compute kernels 
}

\newpage
\subsection{OpenCL -- Open Computing Language}

OpenCL:
\begin{itemize}
    \item \enquote{OpenCL is a trademark of Apple Inc. used under license to the Khronos Group Inc}
    \item working group created 2008~\cite{khronos_opencl_working_group_khronos_2008}
    \item first speciifcation published 2009~\cite{munshi_opencl_2009}
    \item current standard OpenCL 3.0 revision 17 October 24, 2024~\cite{khronos_opencl_working_group_opencl_2024}
    \item Implementations OpenCL 3.0 (\url{https://www.khronos.org/conformance/adopters/conformant-products/opencl})
    \begin{itemize}
        \item Intel
        \item Imagination Technologies
        \item QUALCOMM
        \item Google, Inc. 
        \item Software Freedom Conservancy
        \item Arm Limited
        \item NVIDIA
        \item Software in the Public Interest, Inc.
        \item Verisilicon, Inc.
        \item pocl
    \end{itemize}
    \item many more for OpenCL 2.0 (like AMD, Samsung Electronics, etc.)
\end{itemize}

\begin{itemize}
    \item \acrfull{opencl}
    \item standardized
    \item kernels are simple strings $\rightarrow$ easy to manipulate, tedious to use
    \item kernels must be compiled by yourself using OpenCL provided functions
    \item C API
    \item \gls{jit} compiled
    \item programming model/kernel structure close to \gls{cuda}
\end{itemize}

\begin{advantagebox}
    \begin{itemize}
        \item standardized
        \item supports a wide variety of hardware platforms (many implementations)
    \end{itemize}
\end{advantagebox}

\begin{disadvantagebox}
    \begin{itemize}
        \item kernels must be compiled by hand
        \item error prone
        \item some low-level optimizations not possible
    \end{itemize}
\end{disadvantagebox}


\begin{listing}
    \begin{minted}{cpp}
const char* kernel_source = R"KS(
__kernel void saxpy(const double alpha, __global double *X, __global double *X) {
    const int idx = get_global_id(0);
    Y[idx] = alpha * X[idx] + Y[idx];
}
)KS";

cl_platform_id platform;
clGetPlatformIDs(1, &platform, nullptr);
cl_device_id device;
clGetDeviceIDs(platform, CL_DEVICE_TYPE_GPU, 1, &device, nullptr);
cl_context context = clCreateContext(nullptr, 1, &device, nullptr, nullptr, nullptr);
cl_command_queue queue = clCreateCommandQueue(context, device, 0, nullptr);

cl_mem d_X = clCreateBuffer(context, CL_MEM_READ_ONLY | CL_MEM_COPY_HOST_PTR, N * sizeof(double), X, nullptr);
cl_mem d_Y = clCreateBuffer(context, CL_MEM_READ_ONLY | CL_MEM_COPY_HOST_PTR, N * sizeof(double), Y, nullptr);

cl_program program = clCreateProgramWithSource(context, 1, &kernel_source, nullptr, nullptr);
clBuildProgram(program, 1, &device, nullptr, nullptr, nullptr);
cl_kernel kernel = clCreateKernel(program, "saxpy", nullptr);

clSetKernelArg(kernel, 0, sizeof(double), &alpha);
clSetKernelArg(kernel, 1, sizeof(cl_mem), &d_X);
clSetKernelArg(kernel, 2, sizeof(cl_mem), &d_Y);

size_t global_size = N;
clEnqueueNDRangeKernel(queue, kernel, 1, nullptr, &global_size, nullptr, 0, nullptr, nullptr);
clEnqueueReadBuffer(queue, d_C, CL_TRUE, 0, N * sizeof(double), C, 0, nullptr, nullptr);

clReleaseMemObject(d_X); 
clReleaseMemObject(d_Y);
clReleaseKernel(kernel); 
clReleaseProgram(program);
clReleaseCommandQueue(queue); 
clReleaseContext(context);
    \end{minted}
    \caption{%
    SAXPY example using \gls{opencl}.
    }
    \label{code:opencl_example}
\end{listing}
\newpage
\subsection{SYCL}

SYCL:
\begin{itemize}
    \item first provisional specification March 19, 2014~\cite{khronos_sycl_working_group_khronos_2014}
    \item first official release of SYCL 1.2 specification May 11, 2015~\cite{khronos_sycl_working_group_khronos_2015}
    \item current standard SYCL 2020 Rev 9 August 18, 2024~\cite{khronos_sycl_working_group_sycl_2024}
    \item Implementations (\url{https://www.khronos.org/sycl/#sycl-implementations})
    \begin{itemize}
        \item Intel oneAPI \dpcpp/icpx
        \item AdaptiveCpp
        \item triSYCL
        \item neoSYCL
        \item SimSYCL
    \end{itemize}
\end{itemize}

\begin{itemize}
    \item SYCL
    \item standardized
    \item sinlge-source standard \cpp
    \item \cpp API
    \item easy builtin interoperability with backend specific libraries
    \item different ways to express parallelism
    \item programming model/kernel structure close to \gls{cuda} ($\rightarrow$ automatic \gls{cuda} to SYCL migration tools available)
    \item originally tightly coupled with \gls{opencl}
\end{itemize}


\begin{advantagebox}
    \begin{itemize}
        \item standardized
        \item supports a wide variety of hardware platforms
    \end{itemize}
\end{advantagebox}

\begin{disadvantagebox}
    \begin{itemize}
        \item some low-level optimizations not possible directly in SYCL
    \end{itemize}
\end{disadvantagebox}
\newpage
\subsection{Kokkos}

Kokkos:
\begin{itemize}
    \item first release May 2015~\cite{trott_kokkos_2015}
    \item original Kokkos paper 2024~\cite{carter_edwards_kokkos_2014}
    \item please cite~\cite{trott_kokkos_2022}
    \item current release v 4.5.1 December 23, 2024~\cite{noauthor_release_2024}
\end{itemize}

\begin{itemize}
    \item Kokkos
    \item open-source
    \item \cpp API
    \item interoperability with backend specific libraries
    \item different ways to express parallelism
    \item high-level abstractions, making it easier for new programmers
\end{itemize}


\begin{advantagebox}
    \begin{itemize}
        \item supports a wide variety of hardware platforms with different execution spaces
        \item easy-to-use abstractions like execution and memory spaces
    \end{itemize}
\end{advantagebox}

\begin{disadvantagebox}
    \begin{itemize}
        \item not standardized
        \item abstractions different to \gls{cuda} $\rightarrow$ more difficult to translate from existing \gls{cuda} code
        \item some low-level optimizations not possible directly in Kokkos
    \end{itemize}
\end{disadvantagebox}

\begin{listing}
    \begin{minted}{cpp}
Kokkos::View<double*> d_X("X", N);
Kokkos::View<double*> d_Y("Y", N);
Kokkos::deep_copy(d_X, X);
Kokkos::deep_copy(d_Y, Y);

Kokkos::parallel_for("saxpy", N, KOKKOS_LAMBDA(const int idx) {
    d_Y(idx) = alpha * d_X(idx) + d_Y(idx);
});
Kokkos::fence();

Kokkos::deep_copy(Y, d_Y);
    \end{minted}
    \caption{%
    SAXPY example using Kokkos.
    }
    \label{code:kokkos_example}
\end{listing}

\newpage
\subsection{HIP -- Heterogeneous-computing Interface for Portability}

HIP
\begin{itemize}
    \item first public release 0.82.01 March 08, 2016~\cite{advanced_micro_devices_release_2016}
    \item current release 6.3.3 February, 19 2025~\cite{advanced_micro_devices_release_2025}
    \item current HIP documentation~\cite{advanced_micro_devices_hip_nodate}
    \item HIP for CUDA: \enquote{It’s HIP to be Open}~\cite{advanced_micro_devices_its_2016}
\end{itemize}


\begin{itemize}
    \item \acrfull{hip}
    \item AMD's own vendor-specific framework
    \item C-like API with custom extensions \code{<<<...>>>}
    \item only for AMD and NVIDIA \glspl{gpu}
    \item offload compute kernels from host (i.e., classically \glspl{cpu}) to one or more devices
    \item many vendor-provided, highly optimized libraries available (e.g., hipBLAS, hipSPARSE, hipFFT, $\dots$)
    \item in most cases, conversion between \gls{cuda} and \gls{hip} via simple regex possible
    \item NVIDIA \gls{gpu} support via NVIDIA's compiler stack
\end{itemize}

\begin{advantagebox}
    \begin{itemize}
        \item easy to translate from \gls{cuda}
        \item high optimization potential
    \end{itemize}
\end{advantagebox}

\begin{disadvantagebox}
    \begin{itemize}
        \item not standardized
        \item only AMD and NVIDIA \glspl{gpu}
    \end{itemize}
\end{disadvantagebox}

\begin{listing}
    \begin{minted}{cpp}
__global__ void vector_add(const double alpha, const double *X, double *Y, const int N) {
    const int idx = blockDim.x * blockIdx.x + threadIdx.x;
    if(idx < N) Y[idx] = alpha * X[idx] + Y[idx];
}
    
double *d_X, *d_Y;
hipMalloc(&d_X, N * sizeof(double));
hipMalloc(&d_Y, N * sizeof(double));
hipMemcpy(d_X, X, N * sizeof(double), hipMemcpyHostToDevice);
hipMemcpy(d_Y, Y, N * sizeof(double), hipMemcpyHostToDevice);

vector_add<<<(N + 255) / 256, 256>>>(alpha, d_X, d_Y, N);

hipMemcpy(Y, d_Y, N * sizeof(double), hipMemcpyDeviceToHost);
hipFree(d_X);
hipFree(d_Y);
    \end{minted}
    \caption{%
    SAXPY example using \gls{hip}.
    }
    \label{code:hip_example}
\end{listing}
\newpage
\subsection{Standard \cpp -- stdpar}

standard \cpp
\begin{itemize}
    \item \cpp{17} standard~\cite{international_organization_for_standardization_isoiec_2017}
    \item execution policies introduced with \cpp{17} P0024R2~\cite{hoberock_p0024r2_2016}
    \item \code{std::execution::par_unseq} added in \cpp{20} P1001R2~\cite{meredith_p1001r2_2019}
    \item sender and receiver added in \cpp{26}~\cite{dominiak_p2300r10_2024, baker_p3396r0_2024, niebler_p3325r5_2024}
    \item Implementations:
    \begin{itemize}
        \item GNU GCC + TBB
        \item NVIDIA nvhpc
        \item rocstdpar
        \item AdaptiveCpp
        \item \dpcpp/icpx
    \end{itemize}
\end{itemize}


\begin{itemize}
    \item Standard \cpp
    \item direct part of the \cpp{17} standard
    \item although standard \cpp no all \cpp features allowed inside kernels
    \item parallelization can easily be enabled/disabled using the respective execution policy
\end{itemize}


\begin{advantagebox}
    \begin{itemize}
        \item standardized
        \item supports a wide variety of hardware platforms
        \item extremely easy-to-use if already familiar with \cpp standard algorithms
    \end{itemize}
\end{advantagebox}

\begin{disadvantagebox}
    \begin{itemize}
        \item very high-level, i.e., some optimizations not possible
    \end{itemize}
\end{disadvantagebox}


Other Frameworks:
\begin{itemize}
    \item OpenACC
    \begin{itemize}
        \item first specification November 2011~\cite{openacc_standard_openacc_2011}
        \item first or one of the first OpenACC papers 2012~\cite{hutchison_openacc_2012}
        \item current specification November 2022~\cite{openacc_standard_openacc_2022}
        \item Implementations (\url{https://www.openacc.org/tools})
        \begin{itemize}
            \item HPE Cray Compiler Environment (CCE): OpenACC deprecated since CCE 10.0.0 for C and C++ as of 7 May 2020 (\url{https://support.hpe.com/hpesc/public/docDisplay?docId=a00113907en_us&page=OpenACC_Use.html}, \url{https://support.hpe.com/hpesc/public/docDisplay?docId=a00115296en_us&page=About_the_Cray_Fortran_Reference_Manual.html})
            \item NVIDIA
            \item GCC (\url{https://gcc.gnu.org/wiki/OpenACC})
            \item Omni Compiler Project (\url{https://omni-compiler.org/})
            \item OpenHU (\url{https://github.com/uhhpctools/openuh})
            \item OpenARC (\url{https://github.com/swiftcurrent2018/openarc})
            \item RoseACC (\url{https://github.com/laudarch/RoseACC-workspace})
        \end{itemize}
    \end{itemize}
    
    \item RAJA:
    \begin{itemize}
        \item first release v0.1.0 June 22, 2016~\cite{llnl_release_2016}
        \item first publication 2019~\cite{beckingsale_raja_2019}
        \item current release v2024.07.0 July 24, 2024~\cite{llnl_release_2024}
    \end{itemize}

    \item thrust
    \begin{itemize}
        \item algorithms for NVIDIA GPU and multi-core CPU~\cite{bell_thrust_2012}
        \item first thrust release 1.0.0 June, 2009~\cite{nvidia_corporation_thrust_2009}
        \item current thrust release under the CUDA \cpp Core Libraries 2.8.0 March 03, 2025~\cite{cccl_development_team_release_2023}
    \end{itemize}

    \item Charm++
    \begin{itemize}
        \item GPU Manager library and Hybrid API (HAPI)
        \item source repo~\cite{kale_uiuc-pplcharm_2021}
        \item first release and paper October 1993~\cite{kale_charm_1993}
        \item added support for GPUs~\cite{wesolowski_application_2008}
        \item current release v8.0.0 June 4, 2024~\cite{noauthor_release_2024}
    \end{itemize}

    \item Legion
    \begin{itemize}
        \item \enquote{Legion is a data-centric parallel programming system for writing portable high-performance programs targeted at distributed heterogeneous architectures.}
        \item initial paper including GPU support November, 2012~\cite{bauer_legion_2012}
        \item current release 24.12.0, December 20, 2024~\cite{stanford_release_2024}
    \end{itemize}

    \item Domain-Specific Languages (DSLs)
\end{itemize}

\newpage

\begin{table}[]
    \newcommand{\supported}{\textcolor{ForestGreen}{\faCheckCircle}}
    \newcommand{\partialsupported}{\textcolor{BurntOrange}{\faCheckCircle}}
    \newcommand{\unsupported}{\textcolor{Red}{\faTimesCircle}}
    \centering
    \begin{tblr}{colspec = {Q[l,m] *{6}{X[c,m]}},
                 row{1,2} = {font=\bfseries, bg=gray!50},
                 row{odd[2]} = {bg=gray!25},
                 rowsep=4pt
                 }
        \toprule
                      & \SetCell[c=3]{c}\glspl{cpu} &&             & \SetCell[c=3]{c}\glspl{gpu} &&                            \\
                      & x86\_64      & POWER9       & ARM          & NVIDIA            & AMD               & Intel             \\\midrule
        OpenMP        & \supported   & \supported   & \supported   & \supported        & \supported        & \supported        \\
        CUDA          & \unsupported & \unsupported & \unsupported & \supported        & \unsupported      & \unsupported      \\
        HPX           & \supported   & \supported   & \supported   & \partialsupported & \partialsupported & \partialsupported \\
        OpenCL        & \supported   & \supported   & \supported   & \supported        & \supported        & \supported        \\
        SYCL          & \supported   & \supported   & \supported   & \supported        & \supported        & \supported        \\
        Kokkos        & \supported   & \supported   & \supported   & \supported        & \supported        & \supported        \\
        HIP           & \unsupported & \unsupported & \unsupported & \supported        & \supported        & \unsupported      \\
        Standard \cpp & \supported   & \supported   & \supported   & \supported        & \supported        & \supported        \\
        \bottomrule
    \end{tblr}
    \caption{%
    Either we can use two categories, simply splitting the hardware platforms in \glspl{cpu} and \glspl{gpu}, or split it into more fine-grain categories like different \gls{cpu} and \gls{gpu} architectures. 
    Additional, other hardware platforms inside these categories, e.g., RISC-V \glspl{cpu} or ARM \glspl{gpu} or even completely new categories like \glspl{fpga} or \glspl{tpu} are also possible. 
    For \gls{hpx}, \glspl{gpu} are supported via executors, but the actual compute kernels must be written using another framework like \gls{cuda}, \gls{hip}, SYCL, or Kokkos.
    Note that for the standardized frameworks, it can happen that no single implementation supports all hardware platforms, but only a mixture does. 
    }
    \label{tab:portability_categories}
\end{table}


\begin{table}[]
    \centering
    \begin{tblr}{colspec = {Q[l,m] Q[l,m] Q[l,m] Q[l,m] Q[l,m] Q[l,m]},
                 row{1} = {font=\bfseries, bg=gray!50},
                 row{odd[2]} = {bg=gray!25},
                 rowsep=4pt
                 }
        \toprule
        CUDA             & HIP              & OpenCL            & SYCL              & Kokkos                    & OpenMP \\\midrule
        thread           & thread           & work-item         & work-item         & thread                    & thread \\
        warp             & wavefront        & sub-group         & sub-group         & --                        & -- \\
        {thread\\block}  & {thread\\block}  & work-group        & work-group        & team                      & team \\
        grid             & grid             & nd-range          & nd-range          & league                    & league \\
        register         & register         & {private\\memory} & {private\\memory} & register                  & register \\
        {shared\\memory} & {shared\\memory} & {local\\memory}   & {local\\memory}   & {scratchpad\\memory}      & shared \\
        {global\\memory} & {global\\memory} & {global\\memory}  & {global\\memory}  & \makecell{memory\\spaces} & -- \\
        \bottomrule
    \end{tblr}
    \caption{%
    The different naming conventions for the different portability languages and frameworks. 
    Since my greatest focus lies on SYCL, I will use its naming convention going forward. 
    For Standard \cpp, these notations do not apply since the \cpp standard has no notion of the underlying hardware because it is only defined on an abstract state machine. 
    }
    \label{tab:portability_naming}
\end{table}

